% Generated from locale/tr/content/lambda-calculus/lambda-definability/pairs.tex; do not hand-edit.


\olfileid[tr]{lam}{ldf}{pai}
\olsection{Çiftler ve Öncül}

\begin{defn}
$M$ ile $N$'nin çifti ($\tuple{M,N}$ ile gösterilir) şöyle tanımlanır:
\[
\tuple{M,N} \ident \lambd[f][fMN].
\]
\end{defn}

Sezgisel olarak bu, bir fonksiyon alıp onu çiftin iki elemanına uygulayan bir
fonksiyondur. Bu düşünceye uygun olarak iki terim alıp bunları içeren çifti
döndüren kurucu şöyledir:
\[
  \fn{Pair} \ident \lambd[mn][\lambd[f][fmn]].
\]
Bir çift verildiğinde elemanlarını geri elde etmek isteriz. Bunun için bir
çifti argüman olarak alıp birinci ya da ikinci elemanını döndüren iki erişim
fonksiyonu gerekir:
\begin{align*}
  \fn{Fst} & \ident \lambd[p][p(\lambd[mn][m])]\\
  \fn{Snd} & \ident \lambd[p][p(\lambd[mn][n])].
\end{align*}

\begin{prob}
  $\fn{Fst}$ ve $\fn{Snd}$ erişim fonksiyonlarının neden çalıştığını
  açıklayın.
\end{prob}

Çiftleri kullanarak öncül fonksiyonunu !!{lambda define} edebiliriz:
\[
  \fn{Pred} \ident \lambd[n][\fn{Fst}(n (\lambd[p][\tuple{\fn{Snd}\, {p}, \fn{Succ}(\fn{Snd}\, {p})}]) \tuple{\num 0, \num 0})].
\]
$\num n f x$ teriminin $f^n(x)$'e indirgendiğini hatırlayın. Burada $f$, bir
$p$ çifti alıp $p$'nin ikinci bileşeni ile bu bileşenin ardılını içeren yeni
bir çift döndüren fonksiyondur; $x$ ise $\tuple{\num0,\num0}$ çiftidir.
Dolayısıyla sonuç $n=0$ için $\tuple{\num0,\num0}$, diğer durumlarda
$\tuple{\num{n-1},\num n}$ olur. $\fn{Pred}$ daha sonra sonucun birinci
bileşenini döndürür.

Çıkarma, $\fn{Pred}$ fonksiyonunun $a$'ya $b$ kez uygulanmasıyla tanımlanır:
\[
\fn{Sub} \ident \lambd[ab][b \fn{Pred}\, a].
\]

